%% =====================================================================
%%  CHAPTER 7: WORK PLAN                              (budget 4 pages)
%%  chapters/ch7-plan.tex   (dates synchronised with OrbiTech_14Weeks_WBS_sync.xml)
%% =====================================================================
\pcChapter{plan}{4}{תוכנית עבודה}
\begin{pcguide}[מה כותבים כאן]
תוכנית העבודה מתחילה מהיום שבו התבקש הדוח ומסמנת אבני דרך עד כנס הפרויקטים. בסמסטר ב', בערך במפגש הרביעי או החמישי,
מצפים לראות משהו ראשוני עובד (קריאת נתונים וגרפים ראשונים). בשבוע האחרון מתקיים כנס הפרויקטים, שאליו מוזמנים
ראש המחלקה ונציגי התעשייה האווירית: כל קבוצה מציגה מצגת ומדגימה את הסימולציה. סמנו את \textbf{הנתיב הקריטי}.
תקציב: \latin{1.2} (\latin{WBS}) + \latin{1} (\latin{RACI}) + \latin{1.8} (גאנט).
\end{pcguide}
התוכנית מגדירה פיילוט אקדמי של \latin{14} שבועות בסמסטר ב', שתחילתו ב-\latin{07-03-2027} (יום א', שבוע \latin{1}) וסיומו בכנס הפרויקטים
ב-\latin{10-06-2027}, ולא פריסה ארגונית מלאה. שבוע העבודה הוא ימים א'--ה'; מועדי הסמסטר וכנס הפרויקטים יאושרו בידי
המרצה. שני חברי הצוות הם \StudentTwo\ ו\StudentOne. אב-הטיפוס כולל ליבת \latin{RAG} ותרחיש \latin{3D} מוגבל אחד; תכנון הסצנה, פיתוחה והערכתה נכללים
בסמסטר ואינם נדחים לשלב פריסה עתידי. גורמי התעשייה האווירית והמומחים הם בעלי סמכות האישור התוכני, וסמכותם אינה מוחלפת בסמכות של
סטודנט. \textbf{קודי המשימות \latin{1.1}--\latin{5.2} ואבני הדרך \latin{M0}--\latin{M6} זהים בתרשים ה-\latin{WBS}, במילון החבילות, במטריצת ה-\latin{RACI},
בציר הזמן, בגאנט ובקובץ \latin{MS Project} המעודכן.}

\needspace{4\baselineskip}
\section{פירוק עבודה (\latin{WBS})}
העבודה מפורקת לחמש חבילות ולעשר משימות (איור~\ref{fig:wbs}); לכל חבילה תוצר מוגדר, אחראי ומשך (טבלה~\ref{tab:wbsdict}).
בנוסף הוגדרו שני שבועות רזרבה \latin{(B1, B2)} לפני אבני הדרך המרכזיות \latin{M3} ו-\latin{M6}.
\begin{figure}[!htbp]
\centering
\pcFig{wbs}{%
\pcAutoFit{\begin{pctikz}[x=1cm,y=1cm]
  \node[pcnodeD,text width=4cm] (top) at (0,0) {\tn{פיילוט \latin{RAG} ותרגול \latin{3D} (\latin{14} שבועות)}};
  \foreach \i/\x/\t/\ta/\tb in {1/6.2/{\latin{WP1} אפיון, מדדי בסיס וקורפוס}/{\latin{1.1} מדדים, שאלות בסיס ומדידה ידנית}/{\latin{1.2} קורפוס, אישור שימוש ורובריקה},
                               2/3.1/{\latin{WP2} ארכיטקטורה, קליטה ואינדוקס}/{\latin{2.1} ארכיטקטורה, הרשאות ובידוד רשת}/{\latin{2.2} פענוח, אינדוקס ואחזור},
                               3/0/{\latin{WP3} מנוע הדרכה ותסריט \latin{3D}}/{\latin{3.1} שאלות, תסריט \latin{3D} ומחוון}/{\latin{3.2} מנוע הפקה, מסך אישור וסצנת \latin{3D}},
                               4/-3.1/{\latin{WP4} בדיקות והשוואת \latin{2D/3D}}/{\latin{4.1} הערכת מומחה וניסוי לומדים}/{\latin{4.2} בדיקות אבטחה, תיקונים ורגרסיה},
                               5/-6.2/{\latin{WP5} דוח, הצגה ומסירה}/{\latin{5.1} ניתוח תוצאות ודוח}/{\latin{5.2} הדגמה וכנס הפרויקטים}}{
    \node[pcnode,text width=2.5cm] (w\i) at (\x,-1.9) {\tn{\t}};
    \draw[pcflow,-] (top.south) |- ++(0,-0.55) -| (w\i.north);
    \node[pcleaf,text width=2.5cm,font=\scriptsize\sffamily] (t\i a) at (\x,-3.6) {\tn{\ta}};
    \node[pcleaf,text width=2.5cm,font=\scriptsize\sffamily] (t\i b) at (\x,-4.9) {\tn{\tb}};
    \draw[pcflow,-] (w\i.south) -- (t\i a.north);
    \draw[pcflow,-] (t\i a.south) -- (t\i b.north);
  }
\end{pctikz}}}
\caption{פירוק העבודה (\latin{WBS}) לחמש חבילות עבודה ולעשר משימות}\label{fig:wbs}
\end{figure}

\begin{table}[!htbp]
\centering
\caption{מילון חבילות העבודה: תוצר, אחראי ומשך}\label{tab:wbsdict}
\begin{pctable}{c,H{0.28},H{0.32},C{0.15},C{0.10}}
  \pchead{קוד & חבילת עבודה (משימות) & תוצר מוגדר ואבן דרך & מבצע & שבועות}
  \pcrow{\latin{WP1} & אפיון, מדדי בסיס וקורפוס פומבי (\latin{1.1}, \latin{1.2}) & קורפוס \latin{FAA}, שאלות אמת, רובריקה ומדידת בסיס ידנית (\latin{M1}) & \StudentTwo & \rng{1}{4}}
  \pcrow{\latin{WP2} & ארכיטקטורה מקומית, קליטה ואינדוקס (\latin{2.1}, \latin{2.2}) & תהליך מקומי ללא רשת יוצאת ואחזור ראשון (\latin{M2}) & \StudentOne & \rng{1}{5}}
  \pcrow{\latin{WP3} & מנוע הדרכה, שאלות ותסריט \latin{3D} (\latin{3.1}, \latin{3.2}, רזרבה \latin{B1}) & אב-טיפוס ראשון: אחזור, הפקה, אישור ותצוגת \latin{3D} (\latin{M3}) & \StudentTwo\ ו\StudentOne & \rng{5}{9}}
  \pcrow{\latin{WP4} & בדיקות \latin{RAG}, תרגול \latin{3D} והשוואת \latin{2D} (\latin{4.1}, \latin{4.2}) & ניסוי השוואתי ובדיקות אבטחה ודליפה (\latin{M4}, \latin{M5}) & \StudentTwo\ ו\StudentOne & \rng{10}{12}}
  \pcrow{\latin{WP5} & דוח, הצגה ומסירה (\latin{5.1}, רזרבה \latin{B2}, \latin{5.2}) & דוח, מבחן והדגמת תרחיש \latin{3D} בכנס (\latin{M6}) & \StudentTwo & \rng{12}{14}}
\end{pctable}
\end{table}

\needspace{4\baselineskip}
\section{מטריצת אחריות (\latin{RACI})}
\latin{R} מבצע; \latin{A} האחראי הסופי; \latin{C} מי שמתייעצים עמו; \latin{I} מי שמדווחים לו. לפי מדריך הדוח המכין, בכל שורה בטבלה~\ref{tab:raci}
מופיע \latin{A} אחד בלבד, והוא המנחה האקדמי, ד"ר יורם סגל: הוא האחראי הסופי לאישור כל תוצר מול הקורס. הסטודנטים הם המבצעים
\latin{(R)}. אישורי תוכן ארגוניים (נתונים, סיווג, מקורות) ניתנים בידי בעלי תפקיד מורשים ב-\latin{IAI} בתפקיד \latin{C}, ולא בידי חברי הקבוצה.
\begin{table}[!htbp]
\centering
\caption{מטריצת \latin{RACI} ברמת משימה (קודים זהים ל-\latin{WBS} ולגאנט): בכל שורה \latin{A} אחד --- המנחה האקדמי}\label{tab:raci}
\begin{pctable}{c,H{0.29},C{0.11},C{0.11},C{0.13},C{0.10},C{0.10}}
  \pchead{משימה & תיאור & \StudentOne & \StudentTwo & ד"ר סגל (מנחה) & איש קשר \latin{IAI} & מומחה}
  \pcrow{\latin{1.1} & מדדים, שאלות בסיס ומדידה ידנית & \latin{C} & \latin{R} & \latin{A} & \latin{I} & ---}
  \pcrow{\latin{1.2} & קורפוס, אישור שימוש ורובריקה & \latin{C} & \latin{R} & \latin{A} & \latin{C} & ---}
  \pcrow{\latin{2.1} & ארכיטקטורה, הרשאות ובידוד רשת & \latin{R} & \latin{C} & \latin{A} & \latin{C} & ---}
  \pcrow{\latin{2.2} & פענוח, אינדוקס ואחזור & \latin{R} & \latin{I} & \latin{A} & \latin{I} & ---}
  \pcrow{\latin{3.1} & שאלות, תסריט \latin{3D} ומחוון & \latin{C} & \latin{R} & \latin{A} & \latin{I} & \latin{C}}
  \pcrow{\latin{3.2} & מנוע הפקה, מסך אישור וסצנת \latin{3D} & \latin{R} & \latin{C} & \latin{A} & \latin{I} & ---}
  \pcrow{\latin{4.1} & הערכת מומחה וניסוי לומדים \latin{2D/3D} & \latin{C} & \latin{R} & \latin{A} & \latin{I} & \latin{C}}
  \pcrow{\latin{4.2} & בדיקות אבטחה ודליפה, תיקונים ורגרסיה & \latin{R} & \latin{C} & \latin{A} & \latin{C} & ---}
  \pcrow{\latin{5.1} & ניתוח תוצאות ודוח & \latin{C} & \latin{R} & \latin{A} & \latin{I} & ---}
  \pcrow{\latin{5.2} & הדגמה וכנס הפרויקטים & \latin{R} & \latin{R} & \latin{A} & \latin{I} & ---}
\end{pctable}
\end{table}

\needspace{4\baselineskip}
\section{אבני דרך מהיום ועד כנס הפרויקטים}
איור~\ref{fig:timeline} מציג את הדרך מהגשת הדוח המכין ועד כנס הפרויקטים, ובו שבע אבני הדרך \latin{M0}--\latin{M6} באותם מועדים ובאותם
שבועות כמו בגאנט (איור~\ref{fig:gantt}) ובטבלה~\ref{tab:milestones}. התוצר הראשוני העובד --- אינדקס ואחזור על הקורפוס הפומבי, כלומר קריאת
נתונים ותוצאות ראשונות --- מתוכנן לשבוע \latin{5} \latin{(M2)}, בהתאם לציפייה לראות תוצר ראשוני במפגש הרביעי או החמישי.
\begin{figure}[!htbp]
\centering
\pcFig{timeline}{%
\pcAutoFit{\begin{pctikz}[x=1cm,y=1cm]
  \draw[pcflow,thick,draw=pcDark] (0.5,0) -- (-14.6,0);
  \foreach \x/\lab/\dat/\pos in {0/הגשת הדוח המכין/\latin{25-11-2026}/above, -1.9/\latin{M0} פתיחת הפיילוט/\latin{07-03-2027} (שבוע \latin{1})/below,
      -3.8/\latin{M1} קורפוס ושאלות אמת/\latin{01-04-2027} (שבוע \latin{4})/above, -5.7/\latin{M2} תוצר ראשוני עובד/\latin{08-04-2027} (שבוע \latin{5})/below,
      -7.6/\latin{M3} אב-טיפוס \latin{RAG} ו-\latin{3D}/\latin{06-05-2027} (שבוע \latin{9})/above, -9.5/\latin{M4} ניסוי השוואתי/\latin{20-05-2027} (שבוע \latin{11})/below,
      -11.4/\latin{M5} ולידציה חוזרת/\latin{27-05-2027} (שבוע \latin{12})/above, -13.3/\latin{M6} כנס הפרויקטים/\latin{10-06-2027} (שבוע \latin{14})/below}{
    \fill[pcWarn,draw=pcDark] (\x,0) circle (0.14);
    \node[font=\scriptsize\sffamily,text width=2.6cm,align=center,\pos=4pt] at (\x,0) {\tn{\lab}\\[1pt]\tn{\dat}};
  }
\end{pctikz}}}
\caption{אבני הדרך של תוכנית העבודה, מהגשת הדוח המכין ועד כנס הפרויקטים: מועדים לפי קובץ \latin{MS Project} (תחילת סמסטר \latin{07-03-2027})}\label{fig:timeline}
\figsrc{\latin{OrbiTech\_14Weeks\_WBS\_sync.xml}; מועדי הקורס טרם אושרו בידי המרצה.}
\end{figure}
\begin{table}[!htbp]
\centering
\caption{אבני דרך עד כנס הפרויקטים (מועדים כמו באיורים~\ref{fig:timeline} ו-\ref{fig:gantt})}\label{tab:milestones}
\begin{pctable}{c,H{0.28},C{0.15},C{0.08},H{0.33}}
  \pchead{קוד & אבן דרך & תאריך & שבוע & קריטריון קבלה}
  \pcrow{\latin{M0} & פתיחת הפיילוט & \latin{07-03-2027} & \latin{1} & הגדרת היקף ומקור נתונים}
  \pcrow{\latin{M1} & קורפוס ושאלות אמת (\latin{1.2}) & \latin{01-04-2027} & \latin{4} & מסמכי מקור, גרסאות, רובריקה ומדידת בסיס}
  \pcrow{\latin{M2} & תוצר ראשוני עובד (\latin{2.2}) & \latin{08-04-2027} & \latin{5} & תהליך מקומי ללא רשת יוצאת; אחזור ראשון}
  \pcrow{\latin{M3} & אב-טיפוס ראשון (אחרי \latin{B1}) & \latin{06-05-2027} & \latin{9} & אחזור, הפקה, אישור ותצוגת \latin{3D} בסיסית}
  \pcrow{\latin{M4} & ניסוי השוואתי ראשוני (\latin{4.1}) & \latin{20-05-2027} & \latin{11} & זמן, נאמנות, דיוק חילוץ ורישום החלטות \latin{3D}}
  \pcrow{\latin{M5} & ולידציה חוזרת (\latin{4.2}) & \latin{27-05-2027} & \latin{12} & כל דרישה נבדקה או סומנה ככשל; בדיקת דליפה}
  \pcrow{\latin{M6} & כנס הפרויקטים ומסירה (אחרי \latin{B2}) & \latin{10-06-2027} & \latin{14} & דוח, מבחן והדגמת תרחיש \latin{3D}}
\end{pctable}
\end{table}

\needspace{4\baselineskip}
\section{תרשים גאנט ונתיב קריטי}
התלויות הן מסוג \latin{FS} (סיום-להתחלה) בשתי שרשראות מקבילות. שרשרת התוכן (\StudentTwo):
\latin{1.1} $\leftarrow$ \latin{1.2} $\leftarrow$ \latin{3.1} $\leftarrow$ \latin{B1} $\leftarrow$ \latin{4.1} $\leftarrow$
\latin{5.1} $\leftarrow$ \latin{B2} $\leftarrow$ \latin{5.2}; שרשרת התוכנה (\StudentOne): \latin{2.1} $\leftarrow$
\latin{2.2} $\leftarrow$ \latin{3.2} $\leftarrow$ \latin{B1} $\leftarrow$ \latin{4.2} $\leftarrow$ \latin{B2} $\leftarrow$
\latin{5.2}. בנוסף, \latin{3.2} אינה מתחילה לפני ש-\latin{2.2} מספקת אחזור עובד, ו-\latin{4.1} דורשת את האב-טיפוס של
\latin{M3}. אורך שתי השרשראות \latin{14} שבועות בדיוק, ולכן \textbf{שתיהן על הנתיב הקריטי} (מסגרת אדומה באיור~\ref{fig:gantt}).
שבועות הרזרבה \latin{B1} (שבוע \latin{9}) ו-\latin{B2} (שבוע \latin{13}) הם חלק מהנתיב: הם משמשים לאינטגרציה ולתיקונים,
ואם לא יידרשו --- יוקדם מועד האב-טיפוס או ההגשה.
\begin{figure}[!t]
\centering
\pcFig{gantt}{%
\pcAutoFit{\begin{pctikz}[x=0.55cm,y=0.5cm]
  % week grid (week 1 at the right, time runs right-to-left)
  \foreach \w in {0,...,14}{ \draw[pcGrey!35,thin] ({\pcGsign*\w},0.6) -- ({\pcGsign*\w},-18.6); }
  \foreach \w in {1,...,14}{ \node[font=\scriptsize\sffamily] at ({\pcGsign*(\w-0.5)},1.1) {\w}; }
  \node[font=\scriptsize\sffamily,anchor=west] at (0.25,1.1) {\tn{שבוע}};
  \draw[pcDark,thick] (0,0.6) -- ({\pcGsign*14},0.6);
  % WP1
  \ganttsumS{1}{1}{4}{\latin{WP1} אפיון, מדדי בסיס וקורפוס}
  \ganttbarS{2}{1}{2}{pcPrimary}{נ}{\latin{1.1} מדדים, שאלות בסיס ומדידה ידנית}
  \ganttbarS{3}{3}{4}{pcPrimary}{נ}{\latin{1.2} קורפוס, אישור שימוש ורובריקה}
  % WP2
  \ganttsumS{4}{1}{5}{\latin{WP2} ארכיטקטורה, קליטה ואינדוקס}
  \ganttbarS{5}{1}{3}{pcGood}{א}{\latin{2.1} ארכיטקטורה, הרשאות ובידוד רשת}
  \ganttbarS{6}{4}{5}{pcGood}{א}{\latin{2.2} פענוח, אינדוקס ואחזור}
  % WP3
  \ganttsumS{7}{5}{9}{\latin{WP3} מנוע הדרכה ותסריט \latin{3D}}
  \ganttbarS{8}{5}{8}{pcPrimary}{נ}{\latin{3.1} שאלות, תסריט \latin{3D} ומחוון}
  \ganttbarS{9}{6}{8}{pcGood}{א}{\latin{3.2} מנוע הפקה, מסך אישור וסצנה}
  \ganttbufS{10}{9}{9}{\latin{B1} רזרבה ואינטגרציה לפני \latin{M3}}
  % WP4
  \ganttsumS{11}{10}{12}{\latin{WP4} בדיקות והשוואת \latin{2D/3D}}
  \ganttbarS{12}{10}{11}{pcPrimary}{נ}{\latin{4.1} הערכת מומחה וניסוי לומדים}
  \ganttbarS{13}{10}{12}{pcGood}{א}{\latin{4.2} אבטחה ודליפה, תיקונים ורגרסיה}
  % WP5
  \ganttsumS{14}{12}{14}{\latin{WP5} דוח, הצגה ומסירה}
  \ganttbarS{15}{12}{12}{pcPrimary}{נ}{\latin{5.1} ניתוח תוצאות ודוח}
  \ganttbufS{16}{13}{13}{\latin{B2} רזרבה לפני \latin{M6}}
  \ganttbarS{17}{14}{14}{pcPrimary}{נ}{\latin{5.2} הדגמה וכנס (שניהם)}
  % milestones M0-M6
  \node[anchor=west,text width=4.0cm,align=right,font=\fontsize{7}{8}\selectfont\sffamily\bfseries,inner xsep=1pt] at (0.25,-18.1) {\tn{אבני דרך \latin{M0}--\latin{M6}}};
  \ganttmsS{1}{M0}{18.1}\ganttmsS{4}{M1}{18.1}\ganttmsS{5}{M2}{18.1}\ganttmsS{9}{M3}{18.1}
  \ganttmsS{11}{M4}{18.1}\ganttmsS{12}{M5}{18.1}\ganttmsS{14}{M6}{18.1}
  % legend
  \filldraw[fill=pcPrimary,draw=pcBad,line width=0.7pt] ({\pcGsign*0.2},-19.6) rectangle ({\pcGsign*0.9},-20.2);
  \node[font=\scriptsize\sffamily,anchor=east] at ({\pcGsign*1.1},-19.9) {\tn{\StudentTwo\ (נ)}};
  \filldraw[fill=pcGood,draw=pcBad,line width=0.7pt] ({\pcGsign*4.0},-19.6) rectangle ({\pcGsign*4.7},-20.2);
  \node[font=\scriptsize\sffamily,anchor=east] at ({\pcGsign*4.9},-19.9) {\tn{\StudentOne\ (א)}};
  \filldraw[fill=pcWarn!30,draw=pcWarn!80!black,dashed] ({\pcGsign*7.6},-19.6) rectangle ({\pcGsign*8.3},-20.2);
  \node[font=\scriptsize\sffamily,anchor=east] at ({\pcGsign*8.5},-19.9) {\tn{רזרבה}};
  \draw[pcBad,line width=1.2pt] ({\pcGsign*10.0},-19.6) rectangle ({\pcGsign*10.7},-20.2);
  \node[font=\scriptsize\sffamily,anchor=east] at ({\pcGsign*10.9},-19.9) {\tn{נתיב קריטי}};
  \node[diamond,fill=pcWarn,draw=pcDark,inner sep=1.4pt] at ({\pcGsign*13.2},-19.9) {};
  \node[font=\scriptsize\sffamily,anchor=east] at ({\pcGsign*13.5},-19.9) {\tn{אבן דרך}};
\end{pctikz}}}
\caption{תרשים גאנט ל-\latin{14} שבועות בסמסטר ב': חמש חבילות, עשר משימות ושני שבועות רזרבה; שתי השרשראות על הנתיב הקריטי. אבני הדרך \latin{M0}--\latin{M6} במועדים של טבלה~\ref{tab:milestones} ואיור~\ref{fig:timeline}}\label{fig:gantt}
\figsrc{קובץ \latin{OrbiTech\_14Weeks\_WBS\_sync.xml (MS Project)} המסונכרן עם ה-\latin{WBS}; התרשים מתאר תכנון ולא ביצוע.}
\end{figure}

\textbf{ניהול שני הנתיבים הקריטיים במקביל.} כאשר אין מרווח בשתי השרשראות, עיכוב בכל אחת מהן דוחה את המסירה, ולכן ההתערבות מוגדרת מראש
(טבלה~\ref{tab:paths}). שתי השרשראות נפגשות רק בנקודות האינטגרציה \latin{M3} ו-\latin{M6}; כדי שלא יחסמו זו את זו, הממשק ביניהן
--- סכמת \texttt{\latin{ScenarioConfig}} (פרק~\ref{ch:alt}) --- ננעל בסוף שבוע \latin{5}, וכל שינוי בו מחייב אישור שני הצדדים.
\begin{table}[!htbp]
\centering
\caption{כללי ניהול שתי השרשראות הקריטיות}\label{tab:paths}
\begin{pctable}{H{0.22},H{0.46},H{0.32}}
  \pchead{מנגנון & כלל פעולה & טריגר ופעולה}
  \pcrow{נקודות סנכרון & נפגשים רק ב-\latin{M3} וב-\latin{M6}; ממשק \texttt{\latin{ScenarioConfig}} נעול בשבוע \latin{5} & שינוי בממשק: אישור שני הצדדים ועדכון המסמך}
  \pcrow{מעקב שבועי & בכל מפגש מעודכן סיום מוקדם של כל שרשרת ב-\latin{MS Project}, ומשווים לתכנון & סטייה של יום עבודה אחד: רישום בלבד}
  \pcrow{ניהול רזרבה (\latin{B1}, \latin{B2}) & רזרבה משותפת של \latin{5} ימי עבודה לכל אחת; ירוק: עד \latin{2} ימים; כתום: \latin{3}--\latin{4}; אדום: \latin{5} & אדום: הקפאת הרחבות \latin{3D}, שלוש נקודות החלטה בלבד}
  \pcrow{עדיפות בקונפליקט & כששתי השרשראות דורשות את אותה רזרבה, היא ניתנת לשרשרת שעיכובה חוסם את אבן הדרך הקרובה; בשוויון --- שרשרת התוכנה & החלטה משותפת, מתועדת ביומן שינויים}
  \pcrow{הפרדת משאבים & כל סטודנט מבצע שרשרת אחת ומתייעץ \latin{(C)} בשנייה; עבודה משותפת רק בשבועות האינטגרציה \latin{9} ו-\latin{13} & התלות ההדדית נרשמת במטריצת \latin{RACI}}
\end{pctable}
\end{table}
\textbf{בקרה על שינוי:} בכל מפגש צוות \StudentTwo\ תעדכן חוסרי נתונים, תוכן ומדדים, ו\StudentOne\ יעדכן סטטוס תוכנה, תקלות
ותלויות. שינוי משך משימה ב-\latin{MS Project} מחייב עדכון של פרק זה ושל תרשים הגאנט ואישור המנחה. אם ב-\latin{M3} לא יהיה אב-טיפוס עובד
גם לאחר שבוע הרזרבה, יוקפאו הרחבות ה\threeD{}, ובסצנה הגנרית יישארו שלוש נקודות החלטה בלבד. בכל מקרה אין להתפשר על קישורי
מקור, אישור מומחה, חסימת רשת יוצאת או תיעוד תוצאות.
\pcChapterEnd{plan}
